\section{Introdução}\label{cap:introducao}
% orçamento: 700 palavras

Ptolomeu descreveu o movimento dos planetas com círculos montados sobre círculos. Copérnico mudou o centro e manteve os círculos. Kepler trocou os círculos por elipses, e Newton mostrou de onde as elipses vêm.

Este artigo apresenta um programa que simula o Sistema Solar pela lei da gravitação de Newton, a partir de posições e velocidades do JPL/NASA, e que mostra cada órbita também como uma cadeia de epiciclos. Usamos o programa para responder a uma pergunta antiga com ferramentas modernas: por que os epiciclos funcionavam tão bem, e o que os números de Copérnico têm a ver com a série de Fourier de uma órbita kepleriana.

\subsection{Origem do projeto}

A pergunta vem do trabalho de conclusão de curso de \citeonline{farias2016}, \emph{Copérnico estava certo? Uma análise de epiciclos e séries de Fourier}, defendido na Faculdade UnB Planaltina. O TCC perguntou se as órbitas de Mercúrio, Vênus, Marte, Júpiter e Saturno podem ser escritas como somas de senos e cossenos, e comparou os raios dos epiciclos que Copérnico publicou no \emph{Commentariolus} com os raios obtidos da série de Fourier. Concluiu que dois ou três termos da série já reproduzem bem as órbitas e que os raios de Copérnico têm a ordem de grandeza certa.

Partimos da mesma pergunta com duas diferenças. A órbita vem de uma integração numérica de nove corpos validada contra o JPL, e não de uma elipse ideal. E a série é calculada no tempo, e não no ângulo, o que muda os coeficientes e permite uma comparação direta com o mecanismo de Copérnico (\autoref{sec:param-tempo}).

\subsection{Objetivos}

O objetivo geral é desenvolver um simulador do Sistema Solar fisicamente fundamentado, validado contra dados de referência, e usá-lo para examinar a relação entre a série de Fourier de uma órbita e os modelos de epiciclos. Os objetivos específicos são:

\begin{enumerate}
  \item integrar o problema gravitacional de $N$ corpos com um método que conserve a energia a longo prazo, partindo de condições iniciais do JPL Horizons;
  \item medir o erro das posições calculadas contra a efeméride de referência do JPL;
  \item reconstruir a órbita de cada corpo por série de Fourier complexa e exibir a cadeia de epiciclos correspondente;
  \item comparar, planeta a planeta, os raios do modelo de Copérnico com os coeficientes da série.
\end{enumerate}

\subsection{Contribuições}

\begin{itemize}
  \item um simulador que roda em navegador comum, sem acesso à internet durante o uso, e que publica o próprio erro medido contra o JPL;
  \item uma exposição autocontida da cadeia que vai da lei de Newton à série de Fourier de uma órbita, com a forma fechada dos coeficientes em funções de Bessel conferida numericamente;
  \item a constatação de que a série complexa da órbita kepleriana, até primeira ordem na excentricidade, tem a estrutura do modelo do \emph{Commentariolus} (deslocamento de $\tfrac{3}{2}ae$, deferente e epiciclo de $\tfrac{1}{2}ae$ com frequência dupla), em acordo com análises da literatura \cite{gallavotti2001,sugon2008}, acompanhada de uma comparação numérica com os raios de Copérnico.
\end{itemize}

\subsection{Organização do texto}

A \autoref{cap:historia} resume o contexto histórico, de Ptolomeu a Newton. A \autoref{cap:matematica} desenvolve a série de Fourier de uma órbita kepleriana até a forma fechada dos coeficientes e a aproximação de primeira ordem. A \autoref{sec:metodo} descreve a física do simulador, o integrador e a validação. A \autoref{cap:resultados} reúne os resultados quantitativos, incluindo a comparação com Copérnico. A \autoref{sec:discussao} discute limitações e usos possíveis em sala de aula, e a \autoref{cap:conclusao} conclui. Os apêndices trazem derivações e instruções para reproduzir os resultados.
